\addchap{Abstract}

Axions and axionlike particles (ALPs) are well-motivated dark-matter candidates. Their gradient coupling to nuclear spins produces an oscillating pseudomagnetic field that can be detected using nuclear magnetic resonance (NMR).
In this thesis, signal models, frequency-scanning strategies, and data-analysis methods are developed and applied to proof-of-principle searches with the Cosmic Axion Spin Precession Experiment (CASPEr).

Two dark-matter scenarios are considered: a virialized Milky Way halo, which produces stochastic signals with finite coherence times, and a halo of ALPs gravitationally bound to Earth, which can produce signals with longer coherence times and a potentially enhanced local density.
The frequency-scanning analysis shows that, in the long-measurement limit, coupling sensitivity improves only as $T_\mathrm{tot}^{-1/4}$ with total integration time. Longer nuclear-spin coherence times and frequency steps matched to the sensitive bandwidth are therefore central to efficient searches.

The CASPEr-Gradient-LowField apparatus combines a thermally polarized liquid-methanol sample with SQUID magnetometry.
Searches near \SI{1.35}{\MHz}, corresponding to ALP masses near \SI{5.58}{\neV/c^2}, found no statistically significant candidates.
The Milky Way halo search yields a \SI{95}{\percent} confidence-level upper limit on the ALP-proton coupling of approximately $|g_\mathrm{ap}|\lesssim\SI{1e-2}{\GeV^{-1}}$ over a \SI{240}{\Hz} frequency interval.
The Earth-bound search yields an approximate \SI{95}{\percent} confidence-level limit of $|g_\mathrm{ap}|\lesssim\SI{1e-9}{\GeV^{-1}}$ over a \SI{48}{\Hz} interval, conditional on the assumed Earth-bound halo density and state population.

These measurements demonstrate liquid-state NMR as a method for searching for ALP-proton gradient coupling under distinct halo assumptions.
The very small thermal proton polarization, approximately \num{1.8e-7}, is the principal limitation on sensitivity. Hyperpolarization offers a route to improvements of several orders of magnitude, complemented by lower detector noise and improved field homogeneity.
The signal models, statistical methods, and numerical spin-dynamics simulations developed here provide a basis for future searches with more highly polarized samples.
